% Mathématiques 5e — cours V3
% Opérations sur les nombres

\section{Objectifs}

\begin{itemize}
\item Choisir l'opération adaptée à une situation.
\item Additionner, soustraire, multiplier et diviser pour résoudre des problèmes.
\item Contrôler la vraisemblance d'un résultat.
\item Diviser par un nombre décimal.
\item Enchainer plusieurs opérations.
\item Traduire un problème ou un programme de calcul en une seule expression.
\item Distinguer somme, différence, produit, quotient, termes et facteurs.
\item Utiliser les priorités opératoires.
\item Utiliser la distributivité simple sur des exemples numériques.
\item Utiliser multiples, diviseurs et critères de divisibilité.
\item Mobiliser un algorithme simple de calcul numérique.
\end{itemize}

\section{1. Choisir une opération}

Une opération n'est pas seulement une technique de calcul. Elle traduit une relation entre des quantités.

L'addition permet notamment de réunir des quantités ou de calculer un total.

La soustraction permet de chercher un écart, un reste ou une différence.

La multiplication peut traduire une répétition d'une même quantité ou un agrandissement multiplicatif.

La division peut traduire un partage ou la recherche du nombre de parts.

\begin{exemple}
Une salle comporte $18$ rangées de $24$ sièges.

Le nombre total de sièges se calcule par :
\[
18\times24.
\]

On obtient :
\[
18\times24=432.
\]

La salle comporte donc :
\[
\boxed{432\text{ sièges}}.
\]
\end{exemple}

\section{2. Contrôler la vraisemblance}

Un résultat numérique doit toujours être contrôlé.

On peut utiliser :
\begin{itemize}
\item un ordre de grandeur ;
\item une opération inverse ;
\item une estimation mentale ;
\item une comparaison avec la situation de départ.
\end{itemize}

\begin{exemple}
On calcule :
\[
19{,}8\times5{,}1.
\]

Comme :
\[
20\times5=100,
\]
le résultat exact doit être proche de $100$.

Un résultat comme :
\[
1009{,}8
\]
serait donc manifestement invraisemblable.
\end{exemple}

\section{3. Diviser par un nombre décimal}

En 5e, on apprend à effectuer une division lorsque le diviseur est décimal.

L'idée centrale consiste à transformer la division en une division équivalente avec un diviseur entier.

\begin{propriete}
Multiplier le dividende et le diviseur par un même nombre non nul ne change pas le quotient.
\end{propriete}

Par exemple :
\[
12{,}6\div0{,}3
\]
peut être transformé en :
\[
126\div3.
\]

Donc :
\[
12{,}6\div0{,}3=42.
\]

\section{4. Exemple développé : prix unitaire}

$4{,}5\ \text{kg}$ de fruits coûtent :
\[
13{,}50\ \text{€}.
\]

Le prix pour $1\ \text{kg}$ vaut :
\[
13{,}50\div4{,}5.
\]

On multiplie les deux nombres par $10$ :
\[
13{,}50\div4{,}5
=
135\div45.
\]

Or :
\[
135\div45=3.
\]

Le prix unitaire vaut donc :
\[
\boxed{3\ \text{€ par kg}}.
\]

\section{5. Enchainer plusieurs opérations}

Lorsqu'une expression contient plusieurs opérations, il faut respecter un ordre précis.

\coursefigure{/images/mathematiques/5eme/operations-priorites.svg}{Schéma des priorités opératoires : parenthèses, multiplications et divisions, additions et soustractions.}{Les calculs ne sont pas effectués simplement de gauche à droite : certaines opérations sont prioritaires.}

\begin{propriete}
Dans une expression :
\begin{enumerate}
\item on calcule d'abord les parenthèses ;
\item puis les multiplications et divisions ;
\item enfin les additions et soustractions.
\end{enumerate}
À priorité égale, on calcule de gauche à droite.
\end{propriete}

\section{6. Exemple de priorité}

Considérons :
\[
A=7+3\times5.
\]

On calcule d'abord la multiplication :
\[
3\times5=15.
\]

Puis :
\[
A=7+15=22.
\]

Ainsi :
\[
\boxed{A=22}.
\]

Calculer :
\[
(7+3)\times5
\]
donnerait au contraire :
\[
10\times5=50.
\]

Les parenthèses changent donc le sens du calcul.

\section{7. Nommer un calcul}

Il est important de savoir identifier la nature d'une expression.

Dans :
\[
7+3\times5,
\]
l'opération effectuée en dernier est une addition.

L'expression entière est donc une \textbf{somme}.

Dans :
\[
4\times(6+2),
\]
l'opération effectuée en dernier est une multiplication.

L'expression entière est donc un \textbf{produit}.

\begin{definition}
Dans une somme, les nombres ou expressions additionnés sont les \textbf{termes}.

Dans un produit, les nombres ou expressions multipliés sont les \textbf{facteurs}.
\end{definition}

\section{8. Traduire un problème par une expression}

Une situation peut être résumée par une seule expression numérique.

\begin{exemple}
On achète $3$ cahiers à $2{,}40\ \text{€}$ et un stylo à $1{,}80\ \text{€}$.

Le prix total peut s'écrire :
\[
3\times2{,}40+1{,}80.
\]

On calcule :
\[
3\times2{,}40=7{,}20,
\]
puis :
\[
7{,}20+1{,}80=9.
\]

Le total vaut :
\[
\boxed{9\ \text{€}}.
\]
\end{exemple}

\section{9. Traduire un programme de calcul}

Considérons le programme :
\begin{enumerate}
\item choisir $8$ ;
\item ajouter $5$ ;
\item multiplier le résultat par $3$.
\end{enumerate}

On peut traduire ce programme par :
\[
(8+5)\times3.
\]

Les parenthèses sont indispensables, car le résultat de l'addition doit être multiplié par $3$.

\section{10. Distributivité simple}

La distributivité relie multiplication et addition.

\begin{propriete}
Pour des nombres $k$, $a$ et $b$ :
\[
k(a+b)=ka+kb.
\]

De même :
\[
k(a-b)=ka-kb.
\]
\end{propriete}

Cette propriété peut être utilisée pour calculer plus facilement.

\section{11. Exemple de calcul réfléchi}

Calculons :
\[
17\times99.
\]

On écrit :
\[
99=100-1.
\]

Donc :
\[
17\times99
=
17\times(100-1).
\]

Par distributivité :
\[
17\times100-17\times1.
\]

Ainsi :
\[
1700-17=1683.
\]

Donc :
\[
\boxed{17\times99=1683}.
\]

\section{12. Factoriser numériquement}

La distributivité peut aussi être utilisée dans l'autre sens.

Par exemple :
\[
7\times18+7\times2.
\]

Le facteur commun est :
\[
7.
\]

On écrit :
\[
7\times(18+2).
\]

Donc :
\[
7\times20=140.
\]

Cette écriture est souvent plus rapide.

\section{13. Multiples et diviseurs}

\begin{definition}
Un nombre entier $a$ est un multiple d'un entier $b$ s'il existe un entier $k$ tel que :
\[
a=b\times k.
\]

On dit alors que $b$ est un diviseur de $a$.
\end{definition}

Par exemple :
\[
42=6\times7.
\]

Donc $42$ est un multiple de $6$ et de $7$.

Les nombres $6$ et $7$ sont des diviseurs de $42$.

\section{14. Utiliser les multiples}

Les multiples apparaissent dans de nombreux problèmes.

\begin{exemple}
Un feu clignote toutes les $6$ secondes et un autre toutes les $8$ secondes.

Les multiples de $6$ sont :
\[
6,\ 12,\ 18,\ 24,\ 30,\ldots
\]

Les multiples de $8$ sont :
\[
8,\ 16,\ 24,\ 32,\ldots
\]

Le premier multiple commun positif est :
\[
24.
\]

Les deux feux clignoteront donc ensemble après :
\[
\boxed{24\ \text{s}}.
\]
\end{exemple}

\section{15. Critères de divisibilité}

Certains critères permettent de décider rapidement si un entier est divisible par un autre.

\coursefigure{/images/mathematiques/5eme/operations-divisibilite.svg}{Schéma reliant un nombre entier aux critères de divisibilité par 2, 3, 5, 9 et 10.}{Les critères évitent d'effectuer systématiquement une division.}

\begin{itemize}
\item Divisible par $2$ : le chiffre des unités est pair.
\item Divisible par $5$ : le chiffre des unités est $0$ ou $5$.
\item Divisible par $10$ : le chiffre des unités est $0$.
\item Divisible par $3$ : la somme des chiffres est divisible par $3$.
\item Divisible par $9$ : la somme des chiffres est divisible par $9$.
\end{itemize}

\section{16. Exemple : critère par 3}

Le nombre :
\[
4\,725
\]
a pour somme de chiffres :
\[
4+7+2+5=18.
\]

Comme :
\[
18
\]
est divisible par $3$, le nombre :
\[
4\,725
\]
est divisible par $3$.

Il est aussi divisible par $9$, car :
\[
18=2\times9.
\]

\section{17. Algorithme de divisibilité}

On peut décrire une procédure pour tester la divisibilité par $3$.

\begin{verbatim}
lire le nombre
additionner ses chiffres
si la somme est divisible par 3
    conclure : le nombre est divisible par 3
sinon
    conclure : il ne l'est pas
\end{verbatim}

L'algorithme formalise une méthode répétable.

\section{18. Division euclidienne et divisibilité}

La division euclidienne d'un entier $a$ par un entier $b$ s'écrit :
\[
a=bq+r
\]
avec :
\[
0\le r<b.
\]

Si :
\[
r=0,
\]
alors $a$ est divisible par $b$.

\begin{exemple}
\[
84=7\times12+0.
\]

Donc :
\[
84
\]
est divisible par :
\[
7.
\]
\end{exemple}

\section{19. Priorités avec plusieurs niveaux}

Considérons :
\[
B=18-2\times(5+1).
\]

On calcule la parenthèse :
\[
5+1=6.
\]

Puis la multiplication :
\[
2\times6=12.
\]

Enfin :
\[
18-12=6.
\]

Donc :
\[
\boxed{B=6}.
\]

\section{20. Exemple développé de problème}

Un club achète $12$ ballons à $18{,}50\ \text{€}$ l'unité et bénéficie d'une réduction globale de $25\ \text{€}$.

Le prix avant réduction est :
\[
12\times18{,}50.
\]

On peut écrire le prix final :
\[
12\times18{,}50-25.
\]

Calcul :
\[
12\times18{,}50=222.
\]

Puis :
\[
222-25=197.
\]

Le prix final vaut :
\[
\boxed{197\ \text{€}}.
\]

Un ordre de grandeur donne :
\[
12\times20-25\approx215,
\]
ce qui confirme que $197$ est plausible.

\section{21. Méthode générale}

\begin{methode}
Pour résoudre un problème de calcul numérique :
\begin{enumerate}
\item identifier les grandeurs et les données ;
\item choisir les opérations ;
\item écrire une expression si plusieurs opérations s'enchainent ;
\item respecter les priorités ;
\item effectuer les calculs ;
\item écrire l'unité ;
\item contrôler la vraisemblance du résultat.
\end{enumerate}
\end{methode}

\section{22. Erreurs fréquentes}

\begin{itemize}
\item Calculer une expression uniquement de gauche à droite.
\item Oublier qu'une multiplication est prioritaire sur une addition.
\item Diviser par un décimal sans transformer correctement les deux termes.
\item Confondre multiple et diviseur.
\item Utiliser un critère de divisibilité sur un nombre non entier.
\item Oublier l'unité dans un problème.
\item Donner un résultat numérique sans vérifier s'il est plausible.
\end{itemize}

\section{À retenir}

\begin{aretenir}
Les quatre opérations doivent être choisies en fonction du sens du problème.

Pour un enchainement :
\[
\boxed{\text{parenthèses} \rightarrow \text{multiplications/divisions} \rightarrow \text{additions/soustractions}}.
\]

La distributivité donne :
\[
\boxed{k(a+b)=ka+kb}.
\]

Un multiple s'écrit :
\[
a=b\times k.
\]

Les critères de divisibilité par $2$, $3$, $5$, $9$ et $10$ permettent de tester rapidement de nombreux entiers.

Enfin, tout résultat doit être contrôlé par un ordre de grandeur ou une autre vérification.
\end{aretenir}
